\tikzset{table/.style={
  matrix of nodes,
  row sep=-\pgflinewidth,
  column sep=-\pgflinewidth,
  nodes={rectangle,draw=black,align=left,anchor=center},
  nodes in empty cells,
  every even row/.style={nodes={fill=gray!20}},
  column 1/.style={nodes={text width=19.8em}},
  column 2/.style={nodes={text width=4.8em}},
  column 3/.style={nodes={text width=3.8em}},
  column 4/.style={nodes={text width=19.2em}},
  row 1/.style={nodes={minimum height=1.6em,anchor=center,align=center,fill=darkgray,text=white,font=\bfseries}}}}
\tikzset{tableA/.style={table,
  nodes={minimum height=2.8em},
  column 1/.style={nodes={text width=19.8em}},
  column 2/.style={nodes={text width=4.8em}},
  column 3/.style={nodes={text width=3.8em}},
  column 4/.style={nodes={text width=19.2em}}}}
\tikzset{tableB/.style={table,
  nodes={minimum height=2.8em},
  column 1/.style={nodes={text width=19.8em}},
  column 2/.style={nodes={text width=4.8em}},
  column 3/.style={nodes={text width=6.4em}},
  column 4/.style={nodes={text width=16.6em}}}}
\tikzset{tableModuleName/.style={table,
  nodes={minimum height=1.6em},
  column 1/.style={nodes={text width=16em}},
  column 2/.style={nodes={text width=16em}},
  column 3/.style={nodes={text width=7em}}}}
\tikzset{tableName/.style={table,nodes={minimum height=1.6em}}}
% =============================================================================
\section{Binding Generation}
\label{sec:bind-gen}
% =============================================================================
The most simple form of building an application written in C++ that
depends on \cgal{} consists of the stages below. Let
\myLstinline{\$APP\_SRC\_DIR} point at the root directory of the
application source tree.
\begin{compactenum}
\item Obtain the sources of \cgal{}; let \myLstinline{\$CGAL\_SRC\_DIR}
  point at the root directory of the \cgal{} source tree.
\item Create a build directory for \cgal{}; let \myLstinline{\$CGAL\_DIR}
  point at this build directory.
\item Change directory to \myLstinline{\$CGAL_DIR} and run
  '\myLstinline{cmake [options] \$CGAL\_SRC\_DIR}'.
\item Create a build directory for the application; let
  \myLstinline{\$APP\_DIR} point at this build directory.
\item Change directory to \myLstinline{\$APP\_DIR} and run
  '\myLstinline{cmake [options] \$APP\_SRC\_DIR}' followed by
  '\myLstinline{make}'.
\end{compactenum}

In the procedure above \myLstinline{options} refer to optional
arguments passed to \myLstinline{cmake}, which govern the generation
of the native build environment. When running an application written
in Python that depends on \cgal{} instead of building the application,
the user must build the bindings. Here, \myLstinline{\$APP\_SRC\_DIR}
needs to point at the root directory of the binding source tree,
and \myLstinline{\$APP\_DIR} needs to point at the build directory for
the bindings. The optional arguments passed to \myLstinline{cmake}
in stage~5 mainly consist of variables that determine the objects
to be bound and the unique identifier of each such object, and the
name of the libraries that comprise the bindings, in case the default
library is insufficient.

By default, the name of the generated library
is \myLstinline{CGALPY.so}. However, the user can override the name
with a generated string that maps to the set of bound instances. This
is useful when more then one instance of the same template must be
bound. Each execution of \myLstinline{cmake} followed by make
generates a single library.
% Example: Arrangement_2<Geometry_traits_2, Dcel>
%

The source of \cgal{} consists of several packages.  We use the same
partition and refer to each package as a module. Each module has a
long (and meaningful) name and a short name; see
Table~\ref{tab:module-name} for the list of modules supported so
far. Each module is associated with a \cmake{} Boolean flag that
determines whether to generate bindings for that particular
module. Each module may be associated with zero or more additional
flags that specify selections for bindings within the module. For
example, the long name of the \iiDArrangementsPackage{} module is
\myLstinline{ARRANGEMENT\_2}; its short name is \myLstinline{ARR2};
the Boolean flag that indicates whether to generate bindings for this
module is \myLstinline{CGALPY\_ARRANGEMENT\_2\_BINDINGS}. This module
is associated with the Boolean flag
\myLstinline{CGALPY\_ARR2\_POINT\_LOCATION\_BINDINGS} that indicates
whether to generate bindings for point location queries supported by
the \iiDArrangementsPackage{} package. This module has also a string
flag \myLstinline{CGALPY\_ARR2\_GEOMETRY\_TRAITS\_NAME} that indicates
the geometry traits used to instantiate the
\myLstinline{Arrangement\_2<>} class template.

The Python attribute names of the bound entities of different modules
are gathered in distinct Python scopes to prevent name conflicts. The
name of the scope is the short name of the module in lower-case except
for the first letter (which is in upper-case); for example, the scope
names of the bound types and free functions of
the \iiDArrangementsPackage{} and the \iiiDTriangulationsPackage{}
packages are \myLstinline{Arr2} and
\myLstinline{Tri3}, respectively. Each of these two scopes, for
example, have the attribute \myLstinline{Vertex\_handle} (i.e.,
\myLstinline{Arr2.Vertex\_handle} and
\myLstinline{Tri3.Vertex\_handle}).

As mentioned above, the default name of the generated library can be
overriden. The \cmake{}
flag \myLstinline{CGALPY\_FIXED\_LIBRARY\_NAME} determines whether the
name of the generated library is \myLstinline{CGALPY.so} or not. If
not, it has the prefix \myLstinline{CGALPY_} followed by as many as
substrings as modules, the bindings of which are enabled, separated by
an underscore (\myLstinline{_}). Each such substring starts with the
short name of the module in lower case followed by strings that maps
to the selections for bindings within the module. Each such string is
a single word that starts with a capital letter. For example, the name
\myLstinline{GALPY_kernelEpecInt_Arr2SegPlainPl} of a generated
library, names a library that contains bindings for (i) the
Exact-Predicate-Exact-Construction
(EPEC) \iiDandiiiDGeometryKernelPackage{} module and intersections,
and (ii) the \iiDArrangementsPackage{} module, where
the \myLstinline{Arrangement\_2<>} class template is instantiated with
a traits class that handles segments and a plain DCEL, and point
location queries.

\begin{table}[!]
  \centering\begin{tikzpicture}
    \matrix (first) [tableModuleName] {
Module                  & Name                             & Short Name\\
2D and 3D Kernel        & \myLstinline{KERNEL}             & \myLstinline{KER}\\
2D Arrangements on Surfaces & \myLstinline{ARRANGEMENT\_2} & \myLstinline{ARR2}\\
3D Alpha shapes         & \myLstinline{ALPHA\_SHAPE\_3}    & \myLstinline{AS2}\\
2D Regularized Boolean operations & \myLstinline{BOOLEAN\_SET\_OPERATIONS\_2} & \myLstinline{BSO2}\\
Bounding Volumes        & \myLstinline{BOUNDING\_VOLUMES}  & \myLstinline{BV}\\
2D Convex hulls         & \myLstinline{CONVEX\_HULL\_2}    & \myLstinline{CH2}\\
3D Convex hulls         & \myLstinline{CONVEX\_HULL\_3}    & \myLstinline{CH3}\\
2D Polygons             & \myLstinline{POLYGON\_2}         & \myLstinline{POL2}\\
Polygon partitioning    & \myLstinline{POLYGON\_PARTITIONING} & \myLstinline{PP}\\
2D Minkowski sums       & \myLstinline{MINKOWSKI\_SUM\_2}  & \myLstinline{MS2}\\
Spatial searching       & \myLstinline{SPATIAL\_SEARCHING} & \myLstinline{SS}\\
2D Triangulations       & \myLstinline{TRIANGULATION\_2}   & \myLstinline{TRI2}\\
3D Triangulations       & \myLstinline{TRIANGULATION\_3}   & \myLstinline{TRI3}\\
    };
  \end{tikzpicture}
  \caption{Module names and short names}
  \label{tab:module-name}
\end{table}

The generated library is dynamically linked---it must be so. However,
the library itself can be compiled of either static or dynamic
(dependent) libraries. If you intend to generate and use just a single
library that contains the bindings, you have the freedom to choose
between generating a library compiled of static libraries or dynamic
libraries. However, if you intend to generate and use several
libraries in a single Python module, it is recommended that you use
libraries compiled of dynamic libraries. (Otherwise, different
generated libraries might be compiled of conflicting static
libraries.) The cmake flag \myLstinline{CGALPY\_USE\_SHARED\_LIBS}
indicates whether the generated library is compiled of static or
dynamic libraries; it is true by default.

The content of the library and its name are governed by flags provided
to \myLstinline{cmake}. The number of flags is already large, and it
is expected to grow in the future. All flags have the prefix
\myLstinline{CGALPY_}. In Tables~\ref{tab:args}, \ref{tab:sel},
\ref{tab:ker}, \ref{tab:arr}, \ref{tab:tri}, \ref{tab:as}, and
\ref{tab:ss} this prefix is omitted.

%%%%%%%%
\begin{table}[!htp]
  \tikzset{tableGA/.style={table,
      nodes={minimum height=1.6em},
      row 4/.style={nodes={minimum height=2.8em}},
      column 1/.style={nodes={text width=11em}},
      column 2/.style={nodes={text width=4.8em}},
      column 3/.style={nodes={text width=3.8em}},
      column 4/.style={nodes={text width=28em}}}}
  \centering\begin{tikzpicture}
    \matrix (first) [tableGA] {
Name & Type & Default & Description\\
\myLstinline{USE\_SHARED\_LIBS}    & \myLstinline{Boolean} & \myLstinline{true} & Determines whether to compile shared libraries\\
\myLstinline{BUILD\_SHARED\_LIBS}  & \myLstinline{Boolean} & \myLstinline{true} & Determines whether to generate shared libraries\\
\myLstinline{FIXED\_LIBRARY\_NAME} & \myLstinline{Boolean} & \myLstinline{true} & Determines whether the library name is fixed \myLstinline{cgalpy.so} or set based on other selections\\
    };
  \end{tikzpicture}
  \caption{General Arguments}
  \label{tab:args}
\end{table}

%%%%%%%%
\begin{table}[!]
  \centering\begin{tikzpicture}
    \matrix (first) [tableA] {
Name & Type & Default & Description\\
\myLstinline{ARRANGEMENT\_2\_BINDINGS}              & Boolean & \myLstinline{false} & Determines whether to generate bindings for 2D arrangement instances\\
\myLstinline{ALPHA\_SHAPE\_3\_BINDINGS}             & Boolean & \myLstinline{false} & Determines whether to generate bindings for 3D Alpha shape instances\\
\myLstinline{BOOLEAN\_SET\_OPERATIONS\_2\_BINDINGS} & Boolean & \myLstinline{false} & Determines whether to generate bindings for 2D Boolean set operation instances\\
\myLstinline{BOUNDING\_VOLUMES\_BINDINGS}           & Boolean & \myLstinline{false} & Determines whether to generate bindings for bounding volume instances\\
\myLstinline{CONVEX\_HULL\_2\_BINDINGS}             & Boolean & \myLstinline{false} & Determines whether to generate bindings for 2D convex hull instances\\
\myLstinline{CONVEX\_HULL\_3\_BINDINGS}             & Boolean & \myLstinline{false} & Determines whether to generate bindings for 3D convex hull instances\\
\myLstinline{POLYGON\_2\_BINDINGS}                  & Boolean & \myLstinline{false} & Determines whether to generate bindings for 2D polygon instances\\
\myLstinline{POLYGON\_PARTITIONING\_BINDINGS}       & Boolean & \myLstinline{false} & Determines whether to generate bindings for 2D polygon partitioning instances\\
\myLstinline{MINKOWSKI\_SUM\_2\_BINDINGS}           & Boolean & \myLstinline{false} & Determines whether to generate bindings for 2D Minkowski sum instances\\
\myLstinline{SPATIAL\_SEARCHING\_BINDINGS}          & Boolean & \myLstinline{false} & Determines whether to generate bindings for spatial searching instances\\
\myLstinline{TRIANGULATION\_2\_BINDINGS}            & Boolean & \myLstinline{false} & Determines whether to generate bindings for 2D triangulation instances\\
\myLstinline{TRIANGULATION\_3\_BINDINGS}            & Boolean & \myLstinline{false} & Determines whether to generate bindings for 3D triangulation instances\\
  };
 \end{tikzpicture}
 \caption{Binding Selection}
 \label{tab:sel}
\end{table}

% -----------------------------------------------------------------------------
\subsection{Kernel Bindings}
\label{sec:bind-gen:kernel}
% -----------------------------------------------------------------------------
The \iiDandiiiDGeometryKernelPackage{} package consists of
constant-size non-modifiable geometric primitive objects and
operations on these objects. The objects are sets of points in
$d$-dimensional affine Euclidean space, where $d = 2,3$.  Each point
is uniquely represented either by Cartesian coordinates or by
homogeneous coordinates. An object is represented as a type member of
a kernel type. The kernel type, for example, determines the underlying
number types used for calculations. Table~\ref{tab:ker} lists the \cmake{}
flags associated with the \iiDandiiiDGeometryKernelPackage{} module.
The kernel type is an instance of a chain of C++ class templates. For
convenience, \cgal{} provides the following predefned types of generally
useful kernels:
\begin{compactenum}
\item \myLstinline{Exact_predicates_inexact_constructions_kernel}---provides
  exact geometric predicates, but geometric constructions may be inexact
  due to round-off errors.
  %
\item \myLstinline{Exact_predicates_exact_constructions_kernel}---provides
  exact geometric constructions, in addition to exact geometric predicates.
  %
\item \myLstinline{Exact_predicates_exact_constructions_kernel_with_sqrt}---same
  as\\
  \myLstinline{Exact_predicates_exact_constructions_kernel}, but the
  number type is a model of the concept that requires operations that perform
  square roots, namely \myLstinline{FieldWithSqrt}.\footnote{See \url{https://doc.cgal.org/latest/Algebraic_foundations/classFieldWithSqrt.html}.}
  %
\item \myLstinline{Exact_predicates_exact_constructions_kernel_with_kth_root}---same
  as \\
  \myLstinline{Exact_predicates_exact_constructions_kernel}, but the number
  type is a model of the concept that requires operations that perform $k$-th
  roots, namely
  \myLstinline{FieldWithKthRoot}.\footnote{See \url{https://doc.cgal.org/latest/Algebraic_foundations/classFieldWithKthRoot.html}}
  %
\item \myLstinline{Exact_predicates_exact_constructions_kernel_with_root_of}---same\\
  as \myLstinline{Exact_predicates_exact_constructions_kernel}, but the number
  type is a model of the concept that requires operations that computes the root of univariate polynomial, namely \myLstinline{FieldWithRootOf}.\footnote{See \url{https://doc.cgal.org/latest/Algebraic_foundations/classFieldWithRootOf.html}.}
\end{compactenum}

All the predefined types are of Cartesian
kernels. The \cmake{} flag \myLstinline{KERNEL\_NAME} specifies which
kernel type should be used for the generated bindings.  Currently,
only a limitted predefined types are supported; see Table~\ref{tab:ker}.

%%%%%%%% Kernel
\begin{table}[!htp]
  \centering\begin{tikzpicture}
    \matrix (first) [tableA,row 2/.style={nodes={minimum height=1.6em}}] {
Name & Type & Default & Description\\
\myLstinline{KERNEL\_NAME} & String & \myLstinline{epic} & The kernel type used\\
\myLstinline{KERNEL\_INTERSECTION\_BINDINGS} & Boolean & \myLstinline{true} & Determines whether to generate bindings for intersections\\
  };
 \end{tikzpicture}
 %
 \centering\begin{tikzpicture}
   \matrix (first) [tableName,
     column 1/.style={nodes={text width=8em}},
     column 2/.style={nodes={text width=41em}}
   ] {
   \myLstinline{KERNEL\_NAME} & Predefined Type\\
   \myLstinline{epic} &
     \myLstinline{Exact_predicates_inexact_constructions_kernel}\\
   \myLstinline{epec} &
     \myLstinline{Exact_predicates_exact_constructions_kernel}\\
   \myLstinline{epecws} &
     \myLstinline{Exact_predicates_exact_constructions_kernel_with_sqrt}\\
  };
 \end{tikzpicture}
 \caption{Kernel Arguments}
 \label{tab:ker}
\end{table}

% -----------------------------------------------------------------------------
\subsection{2D Arrangement Bindings}
\label{sec:bind-gen:arr2}
% -----------------------------------------------------------------------------
Given a surface $S$ in \Rthree{} and a set $\calC$ of curves embedded
in this surface, the curves subdivide $S$ into cells of dimension~$0$
(vertices),~$1$ (edges), and~$2$ (faces). This subdivision is the
arrangement $\calA(\calC)$ induced by $\calC$ on $S$. Arrangements
embedded in curved surfaces in \Rthree{} are generalizations of
arrangements embedded in the plane. The \iiDArrangementsPackage{}
package~\cite{cgal:wfzh-a2-21a} can be used to construct, maintain,
alter, and display 2D arrangements embedded in ruled curved
surfaces. It also supports queries on such arrangements, such as point
location. One of the main components of the \iiDArrangementsPackage{}
package is the
\myLstinline{Arrangement_2<Traits,Dcel>} class template. An instance of this
template is used to represent an arrangement embedded in the plane.
Table~\ref{tab:arr} lists the \cmake{} flags associated with
the \iiDArrangementsPackage{} module.  Currently, only instances of
the \myLstinline{Arrangement_2<Traits,Dcel>} class template are
supported. A description of the two template parameters of this class
template follows.

\begin{itemize}
  \item The Traits template-parameter determines the family of planar
  curves that induce the arrangement. The parameter should be
  substituted with a model of the basic arrangement traits concept,
  and optionally additional geometry traits concepts. A model of the
  basic traits concept defines the types of x-monotone curves and
  two-dimensional points and supports basic geometric predicates on
  them. The list of supported traits class template follows.
  %
  \begin{enumerate}
  \item \myLstinline{Arr_non_caching_segment_basic_traits_2<Kernel>}---handles
    segments.
  \item \myLstinline{Arr_segment_traits_2<Kernel>}---handles segments,
    where each segment is represented a by its supporting line in
    addition to its two endpoints.
  \item \myLstinline{Arr_linear_traits_2<Kernel>}---handles linear curves,
    i.e., segments, rays, and lines.
  \item \myLstinline{Arr_polyline_traits_2<SubcurveTraits_2>}
  \item \myLstinline{Arr_polycurve_traits_2<SubcurveTraits_2>}
  \item \myLstinline{Arr_circle_segment_traits_2<Kernel>}
  \item \myLstinline{Arr_conic_traits_2<RatKernel, AlgKernel, NtTraits>}
  \item \myLstinline{Arr_rational_function_traits_2<AlgebraicKernel_d_1>}
  \item \myLstinline{Arr_Bezier_curve_traits_2<RatKernel, AlgKernel, NtTraits>}
  \item \myLstinline{Arr_algebraic_segment_traits_2<Coefficient>}
  \end{enumerate}
  %
\item The Dcel template-parameter should be substituted with a class that
  models the \myLstinline{ArrangementDcel} concept, which is used to
  represent the topological layout of the
  arrangement.\footnote{See \url{https://doc.cgal.org/latest/Arrangement_on_surface_2/classArrangementDcel.html}.}
  This parameter is substituted
  with \myLstinline{Arr_default_dcel<Traits>} by default. In many
  applications it is necessary to extend the DCEL features. This is
  done by substituting the Dcel parameter with a different type; see
  Section Extending the DCEL for further explanations and examples.
\end{itemize}

%%%%%%%% Arrangement
\begin{table}[!htp]
  \centering\begin{tikzpicture}
  \matrix (first) [tableB,
    row 3/.style={nodes={minimum height=1.6em}},
    row 4/.style={nodes={minimum height=4.0em}}
  ] {
    Name & Type & Default & Description\\
    \myLstinline{ARR2\_GEOMETRY\_TRAITS\_NAME}    & String  &
      \myLstinline{segment}     & The geometry traits of a 2D arrangement\\
    \myLstinline{ARR2\_DCEL\_NAME}                & String  &
      \myLstinline{allExtended} & The DCEL of a 2D arrangement\\
    \myLstinline{ARR2\_POINT\_LOCATION\_BINDINGS} & Boolean &
      \myLstinline{true}        & Determines whether to generate bindings for point location and vertical ray shooting queries\\
  };
 \end{tikzpicture}
 %
 \centering\begin{tikzpicture}
   \matrix (first) [tableName,
     column 1/.style={nodes={text width=18em}},
     column 2/.style={nodes={text width=31em}}
   ] {
     \myLstinline{ARR2\_GEOMETRY\_TRAITS\_NAME} & Type\\
     \myLstinline{nonCachingSegment} &
       \myLstinline{Arr_non_caching_segment_basic_traits_2<Kernel>}\\
     \myLstinline{segment} & \myLstinline{Arr_segment_traits_2<Kernel>}\\
     \myLstinline{linear} & \myLstinline{Arr_linear_traits_2<Kernel>}\\
     \myLstinline{conic} &
       \myLstinline{Arr_conic_traits_2<RatKernel, AlgKernel, NtTraits>}\\
     \myLstinline{circleSegment} &
       \myLstinline{Arr_circle_segment_traits_2<Kernel>}\\
     \myLstinline{algebraic} &
       \myLstinline{Arr_algebraic_segment_traits_2<Coefficient>}\\
   };
 \end{tikzpicture}
 %
 \centering\begin{tikzpicture}
   \matrix (first) [tableName,
     column 1/.style={nodes={text width=12em}},
     column 2/.style={nodes={text width=37em}}
   ] {
     \myLstinline{ARR2\_DCEL\_NAME} & Type\\
     \myLstinline{plain} & \myLstinline{Arr_default_dcel<GeometryTraits>}\\
     \myLstinline{vertexExtended} &
       \myLstinline{Arr_vertex_extended_dcel<Traits, Data>}\\
     \myLstinline{halfedgeExtended} &
       \myLstinline{Arr_halfedge_extended_dcel<Traits, Data>}\\
     \myLstinline{faceExtended} &
       \myLstinline{Arr_face_extended_dcel<Traits, Data>}\\
     \myLstinline{allExtended} &
       \myLstinline{Arr_extended_dcel<Traits, VertexData, HalfedgeData, FaceData>}\\
  };
 \end{tikzpicture}
 \caption{2D Arrangement Arguments}
 \label{tab:arr}
\end{table}

% -----------------------------------------------------------------------------
\subsection{3D Triangulation Bindings}
\label{sec:bind-gen:tri3}
% -----------------------------------------------------------------------------

%%%%%%%% 3D Triangulation
\begin{table}[!htp]
  \centering\begin{tikzpicture}
    \matrix (first) [tableB,
       row 2/.style={nodes={minimum height=1.6em}},
       row 3/.style={nodes={minimum height=1.6em}},
       row 4/.style={nodes={minimum height=1.6em}},
       row 6/.style={nodes={minimum height=1.6em}},
       row 7/.style={nodes={minimum height=1.6em}}
    ] {
Name & Type &  Default & Description\\
\myLstinline{TRI3\_NAME}                   & String & \myLstinline{regular}    & The 3D triangulation type\\
\myLstinline{TRI3\_VERTEX\_BASE\_NAME}     & String & \myLstinline{plain}      & \\
\myLstinline{TRI3\_CELL\_BASE\_NAME}       & String & \myLstinline{plain}      & \\
\myLstinline{TRI3\_TRAITS\_NAME}           & String & \myLstinline{kernel}     & The geometry traits of a 2D triangulation\\
\myLstinline{TRI3\_CONCURRENCY\_NAME}      & String & \myLstinline{sequential} & The concurrency method\\
\myLstinline{TRI3\_LOCATION\_POLICY\_NAME} & String & \myLstinline{compact}    & The location policy\\
  };
  \end{tikzpicture}
  %
  \centering\begin{tikzpicture}
    \matrix (first) [tableName,
      column 1/.style={nodes={text width=12em}},
      column 2/.style={nodes={text width=37em}}
    ] {
    \myLstinline{TRI3\_NAME} & Type\\
    \myLstinline{plain} & \\
    \myLstinline{regular} & \\
    \myLstinline{delaunay} & \\
    \myLstinline{periodic3Delaunay} & \\
  };
  \end{tikzpicture}
  %
  \centering\begin{tikzpicture}
    \matrix (first) [tableName,
      column 1/.style={nodes={text width=19em}},
      column 2/.style={nodes={text width=30em}}
    ] {
    \myLstinline{TRI3\_VERTEX|CELL\_BASE\_NAME} & Type\\
    \myLstinline{plain} & \\
    \myLstinline{plainWithInfo} & \\
    \myLstinline{regular} & \\
    \myLstinline{regualrWithInfo} & \\
    \myLstinline{alphaShape} & \\
    \myLstinline{alphaShapeWithInfo} & \\
    \myLstinline{alphaShapeRegular} & \\
    \myLstinline{alphaShapeRegularWithInfo} & \\
    \myLstinline{fixedAlphaShape} & \\
    \myLstinline{fixedAlphaShapeWithInfo} & \\
    \myLstinline{fixedAlphaShapeRegular} & \\
    \myLstinline{fixedAlphaShapeRegularWithInfo} & \\
  };
  \end{tikzpicture}
  %
  \centering\begin{tikzpicture}
    \matrix (first) [tableName,
      column 1/.style={nodes={text width=12em}},
      column 2/.style={nodes={text width=37em}}
    ] {
    \myLstinline{TRI3\_TRAITS\_NAME} & Type\\
    \myLstinline{kernel} & \\
    \myLstinline{periodic3Delaunay} & \\
  };
  \end{tikzpicture}
  %
  \centering\begin{tikzpicture}
    \matrix (first) [tableName,
      column 1/.style={nodes={text width=15em}},
      column 2/.style={nodes={text width=34em}}
    ] {
    \myLstinline{TRI3\_CONCURRENCY\_NAME} & Type\\
    \myLstinline{sequential} & \\
    \myLstinline{parallel} & \\
  };
  \end{tikzpicture}
  %
  \centering\begin{tikzpicture}
    \matrix (first) [tableName,
      column 1/.style={nodes={text width=18em}},
      column 2/.style={nodes={text width=31em}}
    ] {
    \myLstinline{TRI3\_LOCATION\_POLICY\_NAME} & Type\\
    \myLstinline{fast} & \\
    \myLstinline{compact} & \\
  };
  \end{tikzpicture}
  \caption{3D Triangulation Arguments}
  \label{tab:tri}
\end{table}

%%%%%%%%
\begin{table}[!htp]
  \centering\begin{tikzpicture}
    \matrix (first) [tableA] {
Name & Type &  Default & Description\\
\myLstinline{AS3\_NAME}              & String  & \myLstinline{plain} & The 3D Alpha shape type, either plain or fixed\\
\myLstinline{AS3\_EXACT\_COMPARISON} & Boolean & \myLstinline{false} & Determines whether to apply exact comparisons\\
  };
 \end{tikzpicture}
 \caption{3D Alpha Shape Arguments}
 \label{tab:as}
\end{table}

%%%%%%%%
\begin{table}[!htp]
  \centering\begin{tikzpicture}
    \matrix (first) [tableA] {
Name & Type &  Default & Description\\
\myLstinline{SPATIAL\_SEARCHING\_DIMENSION} & Integer & \myLstinline{2} & The dimension of spatial searching related classes\\
  };
 \end{tikzpicture}
 \caption{Spatial Searching Arguments}
 \label{tab:ss}
\end{table}
